\chapter{Experimental Setup}
\label{ch:experiments}

\section{Dataset}

The primary corpus is Metamon Gen 9 OU human replays. The scale target is at
least fifty thousand battles after filtering corrupted or extremely short games.
Observations use first-person reconstructed from spectator logs. Actions come
from human moves. Belief labels come from full-state opponent annotations.

Random-bot logs from poke-env support pipeline validation before human data is
available at full scale.

\section{Data Splits}

We report three complementary split strategies.

\subsection{Split A: Random Battle Split}

Standard 80/10/10 train, validation, and test partition by battle ID. This
measures in-distribution imitation and belief accuracy.

\subsection{Split B: Held-Out Team Cores}

Each opponent team receives a canonical key from sorted species names. Teams in
a designated held-out set never appear in training. Test battles use opponents
whose team key was unseen during training.

\subsection{Split C: Held-Out Set Combinations}

We hold out specific tuples of species, item, ability, and move quartet. This
tests fine-grained combinatorial generalization beyond species identity alone.

\subsection{Split D: Temporal Split (Optional)}

When metadata permits, train on an earlier ladder period and test on a later
period to measure temporal drift.

\section{Belief Metrics}

For each hidden-attribute type and model we report
\begin{itemize}
  \item top-1 accuracy on masked unrevealed slots,
  \item top-3 accuracy where label space is large,
  \item expected calibration error (ECE),
  \item Brier score for multi-class predictions,
  \item reveal-turn anticipation (turns before official reveal when top-1 becomes
    correct).
\end{itemize}

We plot belief accuracy versus turn number to show uncertainty narrowing over
time.

\section{Policy Metrics}

Policy evaluation uses
\begin{itemize}
  \item BC cross-entropy and top-1 action match versus human actions,
  \item offline policy value estimates when Q-learning is used,
  \item live win rate on poke-env against fixed opponents,
  \item generalization gap (in-distribution metric minus out-of-distribution
    metric).
\end{itemize}

The generalization gap is the primary summary statistic for RQ3.

\section{Ablations}

\begin{table}[ht]
  \centering
  \caption{Planned ablation studies.}
  \label{tab:ablations}
  \begin{tabular}{ll}
    \toprule
    Ablation & What it tests \\
    \midrule
    Context $N \in \{4,8,16,32\}$ & History length for belief \\
    Remove within-turn attention & Value of 13-token structure \\
    Bidirectional turn encoder & Future leakage effects \\
    $\lambda_\mathrm{bel} = 0$ vs tuned & Necessity of belief loss \\
    Belief fusion vs shared trunk & Explicit belief features \\
    CQL vs AMAGO vs BC-only & Offline RL contribution \\
    \bottomrule
  \end{tabular}
\end{table}

\section{Statistical Reporting}

We report mean $\pm$ 95\% bootstrap confidence intervals over battles. Paired
comparisons use the same test battles for each model. Win-rate significance
uses paired bootstrap or McNemar tests where appropriate.

\section{Minimum Viable Result}

The thesis meets its pass bar if at least two of the following hold on
combinatorial splits.
\begin{enumerate}
  \item Belief-augmented TurnTransformer beats classical and recurrent belief
    baselines on aggregate hidden-species prediction.
  \item Belief-augmented BC beats BC-only TurnTransformer on OOD team metrics
    with a smaller generalization gap.
  \item Offline RL fine-tuning improves live win rate over BC-only without
    search at inference.
  \item Ablations confirm causal history modeling contributes beyond flat
    sequence baselines.
\end{enumerate}

\section{Compute and Reproducibility}

Training runs on a single GPU when available with CPU fallback. Checkpoints,
vocabularies, and split manifests are saved under \texttt{checkpoints/} and
\texttt{logs/}. The thesis PDF is built from \texttt{thesis/main.tex}.
